\section{国际知识产权法律制度}

\begin{enumerate}
  \item 保护工业产权（\zd{专利和商标}）巴黎公约。四个基本原则。
        \begin{enumerate}
          \item 国民待遇原则。例外。获取程序方面。在我国外国人必须要找专代。
          \item 优先权原则。发明和实用新型 12 月。外观设计和商标 6 月。首次申请日就叫做\zd{优先权日}。
          \item 临时性保护原则。国家举办的国际展览会。首次申请日提前到公开展出的第 1 天。
          \item \zd{独立性原则}。
        \end{enumerate}
  \item 保护文学艺术作品伯尔尼公约（\zd{著作权}）。
        \begin{enumerate}
          \item \zd{双国籍国民待遇}。考的最多。作者国籍和作品国籍。
          \item 自动保护原则。\zd{创作完成自动保护}。
          \item \zd{独立性原则}。以\zd{本国}所规定的\zd{水平}保护，不违反公约最低水平即可，高高低低随便。
        \end{enumerate}
  \item TRIPS 协议/与贸易有关的知识产权协定。WTO 的一个知识产权协议。更高标准更严要求。
        \begin{enumerate}
          \item 首次将\zd{最惠国待遇原则}引入知识产权的国际保护领域。WTO 每一个协议中都有一个最惠国待遇条款。
          \item \zd{更高}水平的立法保护。
          \item \zd{更严格}的执法手段，民事、行政和刑事程序。\zd{必须包括三套}。
          \item \zd{凡是符合中国知识产权法的就一定符合 TRIPS 协议}。
        \end{enumerate}
\end{enumerate}
